\setcounter{section}{2}

\Needspace{5\baselineskip}
\hypertarget{ux591aux667aux80fdux4f53ux53caux5176ux8badux7ec3}{%
\section{多智能体及其训练}\label{ux591aux667aux80fdux4f53ux53caux5176ux8badux7ec3}}

\Needspace{5\baselineskip}
\hypertarget{ux4e00ux80ccux666fux4e0eux4f18ux52bf}{%
\subsection{一、背景与优势}\label{ux4e00ux80ccux666fux4e0eux4f18ux52bf}}

\begin{enumerate}
\def\labelenumi{\arabic{enumi}.}
\item
  单智能体的局限性：上下文窗口有限；推理深度和复杂处理的局限性；多步内部不透明，可解释性弱；垂直专业领域知识不足；容易出现``幻觉''，自我纠错能力弱；知识``冻结''在训练时间节点。
\item
  多智能体系统的优势：避免上下文过长（如可以把代码运行结果先给另一个Agent处理，以控制写代码的Agent上下文长度）；专业化分工（类似于人类社会的分工协作）；抗风险能力强，单个节点故障影响低；多个Agent可以并行处理；灵活性强；可以互相纠错，降低``幻觉''率；Agent间交互透明，可解释性强。
\end{enumerate}

\Needspace{5\baselineskip}
\hypertarget{ux4e8cux591aux667aux80fdux4f53ux534fux4f5cux7684ux5e38ux89c1ux6a21ux5f0f}{%
\subsection{二、多智能体协作的常见模式}\label{ux4e8cux591aux667aux80fdux4f53ux534fux4f5cux7684ux5e38ux89c1ux6a21ux5f0f}}

\begin{enumerate}
\def\labelenumi{\arabic{enumi}.}
\item
  串行式：多步骤推理任务拆解，如提出问题-设计实验-数据分析-论文撰写各一个Agent。
\item
  角色式：扮演不同专家角色，从不同角度分析问题，协商合作，如产品经历、工程师等。
\item
  并行式：如大规模数据集分割开来，不同Agent分别处理。
\item
  辩论式：正方、反方Agent相互辩论、``找漏洞''、质疑，另有负责裁判和检验事实的Agent。这也是避免模型``幻觉''导致误导的重要方法。
\item
  融合式：不同Agent从不同角度分析同一问题，再用简单投票、按过去准确度加权投票、输出融合等方式得出答案，减少单一模型的偏见和错误。
\item
  外部交互式：不同专门的Agent负责与不同的工具、数据库等外部环境交互，如分别负责信息检索、数学计算、文档编写等。
\item
  计划-执行-检验式：一个Agent负责制定计划，一个Agent执行，最后还有Agent独立检验以保证输出符合要求，避免``既当运动员，又当裁判员''。在容错率低的情况下，还可以引入多个串行的检验Agent等。
\end{enumerate}

以上模式也可以根据需要相结合。

\Needspace{5\baselineskip}
\hypertarget{ux4e09ux591aux667aux80fdux4f53ux7684ux5178ux578bux67b6ux6784}{%
\subsection{三、多智能体的典型架构}\label{ux4e09ux591aux667aux80fdux4f53ux7684ux5178ux578bux67b6ux6784}}

\begin{enumerate}
\def\labelenumi{\arabic{enumi}.}
\item
  网络型：每个智能体都可以和其他智能体通信。
\item
  监督型：每个智能体只和一个监督智能体通信。
\item
  层级型：对监督型的扩展，每个智能体和它的``直接上级''智能体通信。
\end{enumerate}

\Needspace{5\baselineskip}
\hypertarget{ux56dbux591aux667aux80fdux4f53ux4e00ux5b9aux6bd4ux5355ux667aux80fdux4f53ux597dux5417}{%
\subsection{四、多智能体一定比单智能体好吗？}\label{ux56dbux591aux667aux80fdux4f53ux4e00ux5b9aux6bd4ux5355ux667aux80fdux4f53ux597dux5417}}

\Needspace{5\baselineskip}
如果不做好上下文工程，盲目使用``多智能体并行架构''可能导致灾难性的后果。如一个开发Flappy Bird游戏的案例：

设想：为了加快速度，让Agent把任务拆解，分给Subagent 1和Subagent 2并行去做。

Subagent 1负责背景：因为它不知道另一边在做什么，结果做了一个超级马里奥风格的背景。

Subagent 2负责鸟：它做了一只不像游戏里的鸟，移动方式也不对。

结果：两个做出来的东西根本拼不到一起。

尝试修复：为了解决割裂，我们尝试把原始任务作为上下文复制给每个子智能体。

依然失败：现实中的任务有很多``隐性知识''是在对话和工具调用中产生的（比如用户中途说``把鸟改成红色的''）。如果子智能体看不到这些动态产生的中间状态，依然会误解任务。

最终瓶颈（上下文溢出）：如果你试图把所有发生过的事情都塞给每一个子智能体，对于复杂任务，Token消耗会瞬间爆炸，导致上下文窗口溢出。

总之，代码生成这种任务需要极高的逻辑连贯性（高耦合）。前面的变量定义影响后面的函数调用，上下文必须高度一致。如果像做阅读理解那样切分任务，代码就跑不通了。多智能体是否优于单智能体需要根据任务而定，架构是服务于任务的。如果一定要用多智能体，则需要探索更好地增强不同Agent之间上下文的协同性的方法。

\Needspace{5\baselineskip}
\hypertarget{ux4e94ux5927ux6a21ux578bux5e7bux89c9ux95eeux9898ux7684ux591aux667aux80fdux4f53ux89e3ux51b3ux65b9ux6848}{%
\subsection{五、大模型``幻觉''问题的多智能体解决方案}\label{ux4e94ux5927ux6a21ux578bux5e7bux89c9ux95eeux9898ux7684ux591aux667aux80fdux4f53ux89e3ux51b3ux65b9ux6848}}

通过多智能体（Multi-Agent）协作、互评和辩论来降低幻觉，是当前大模型研究中的常见技术路径之一。

\Needspace{4\baselineskip}
\begin{enumerate}
\def\labelenumi{\arabic{enumi}.}
\tightlist
\item
  数学原理
\end{enumerate}

假设我们有 \(n\) 个相互独立的模型（智能体），每个模型针对同一个问题犯``低级错误''的概率为 \(p\)（其中 \(0<p<0.5\)，即模型正确的概率大于出错的概率）。

\Needspace{5\baselineskip}
如果采用``多数服从少数''或者``全票通过''的机制，可以推导错误率的变化。情境 A：要求所有模型一致才输出（极为严格的校验）。多个模型同时犯下同一个特定错误的概率 \(P_{\mathrm{error}}\) 为：

\[
P_{\mathrm{error}}=\prod_{i=1}^{n}P(E_i)=p^n
\]

当 \(n\) 增加时，由于 \(p<1\)，\(p^n\) 会以指数级速度趋近于 0。

\Needspace{5\baselineskip}
情境 B：多数投票机制（Majority Voting）。假设有 \(n\) 个模型（\(n\) 为奇数），只要超过半数的模型（\(k>n/2\)）给出正确答案，系统就输出正确结果。根据二项分布，系统最终正确的概率 \(P_{\mathrm{correct\_sys}}\) 为：

\[
P_{\mathrm{correct\_sys}}=\sum_{k=\lfloor n/2\rfloor+1}^{n}\binom{n}{k}(1-p)^k p^{n-k}
\]

根据孔多塞定理，只要单个个体的正确率 \((1-p)>0.5\)，随着 \(n\to\infty\)，系统整体正确的概率 \(P_{\mathrm{correct\_sys}}\to 1\)。

\Needspace{4\baselineskip}
\begin{enumerate}
\def\labelenumi{\arabic{enumi}.}
\setcounter{enumi}{1}
\tightlist
\item
  现实挑战：独立性假设往往并不成立
\end{enumerate}

\Needspace{8\baselineskip}
（1）训练数据的同源性，导致不同模型输出的相关性强；

\Needspace{8\baselineskip}
（2）RLHF让模型倾向于迎合其他模型的回答；

（3）对于某些逻辑陷阱或极其专业的知识，验证可能比生成更难。

\Needspace{4\baselineskip}
\begin{enumerate}
\def\labelenumi{\arabic{enumi}.}
\setcounter{enumi}{2}
\tightlist
\item
  目前的技术路线
\end{enumerate}

\Needspace{8\baselineskip}
（1）思维链的自我一致性；

\Needspace{8\baselineskip}
（2）不同智能体生成后多数投票；

\Needspace{8\baselineskip}
（3）使用完全不同架构或不同侧重点的模型进行互评；

\Needspace{8\baselineskip}
（4）让多智能体扮演不同角色，进行对抗性辩论；

（5）引入外部工具（如代码解释器）进行验证。

其中，随机性错误（低级计算错误、口误）通过（1）（2）（3）即可有效解决；系统性错误（知识盲区、逻辑陷阱）需要通过（4）（5）才能解决。

\Needspace{5\baselineskip}
\hypertarget{ux516dux591aux667aux80fdux4f53ux5e76ux884cux5f3aux5316ux5b66ux4e60parallel-agent-reinforcement-learning-parl}{%
\subsection{六、多智能体并行强化学习（Parallel-Agent Reinforcement Learning, PARL）}\label{ux516dux591aux667aux80fdux4f53ux5e76ux884cux5f3aux5316ux5b66ux4e60parallel-agent-reinforcement-learning-parl}}

\Needspace{4\baselineskip}
\begin{enumerate}
\def\labelenumi{\arabic{enumi}.}
\tightlist
\item
  核心架构
\end{enumerate}

以Kimi K2.5 Agent Swarm为例，推理过程中，当输入一个任务时，一个编排器负责确定输入任务需要子智能体的数量及角色（以上下文形式提供给智能体）、分配任务、调度智能体；每个子智能体都是同一个大模型的``分身''，权重完全相同。

\Needspace{4\baselineskip}
\begin{enumerate}
\def\labelenumi{\arabic{enumi}.}
\setcounter{enumi}{1}
\tightlist
\item
  训练方法
\end{enumerate}

先训练子智能体（大模型）；在子智能体训练完毕后，冻结子智能体的参数，用端到端的强化学习方法训练编排器（刚开始先用小模型充当子智能体，后续再换为大模型）。

\Needspace{4\baselineskip}
\begin{enumerate}
\def\labelenumi{\arabic{enumi}.}
\setcounter{enumi}{2}
\tightlist
\item
  奖励函数
\end{enumerate}

\Needspace{5\baselineskip}
训练信号由三部分组成，通过强化学习优化编排器的并行调度策略：

\[
r_{\mathrm{PARL}}(x,y)=\lambda_1\cdot\underbrace{r_{\mathrm{parallel}}}_{\mathcjk{实例化奖励}}+\lambda_2\cdot\underbrace{r_{\mathrm{finish}}}_{\mathcjk{完成率奖励}}+\underbrace{r_{\mathrm{perf}}(x,y)}_{\mathcjk{任务结果奖励}}
\]

\Needspace{8\baselineskip}
（1）并行奖励 r\_parallel：缓解``串行崩溃''

问题：编排器可能陷入局部最优，默认使用单智能体串行执行。

机制：通过奖励子智能体的实例化行为，鼓励探索并发调度空间。

作用：强制系统尝试并行化解构，而非顺序执行。

\Needspace{8\baselineskip}
（2）完成奖励 r\_finish：防止``虚假并行''

问题：编排器可能通过生成大量无意义的子智能体，增加并行指标但不实际分解任务。

机制：仅对成功完成（如成功返回查询结果、成功输出最终答案、成功解析网页）的子任务给予奖励。

作用：确保任务分解的可行性，引导策略生成有效的并行结构。

\Needspace{8\baselineskip}
（3）性能奖励 r\_perf：保证最终质量

基于任务最终结果好坏的可验证奖励（如答案正确性）。

随着训练进行，λ1和 λ2逐渐退火至0，使策略最终优化主要目标。

对于网页设计审美、开放式生成等无法通过规则或二元对错验证的任务，可采用GRMs（生成式奖励模型）提供细粒度评估。GRMs可能有一个（不同任务用不同提示词）或多个，设计提示词使之采取多维度综合连续评分。最终得分由可验证奖励（如搜索到的事实是否正确）+GRMs评分加权得到，尽可能避免过度Reward Hacking。

\Needspace{4\baselineskip}
\begin{enumerate}
\def\labelenumi{\arabic{enumi}.}
\setcounter{enumi}{3}
\tightlist
\item
  任务设计
\end{enumerate}

PARL的训练任务专门设计用于压力测试顺序执行的极限，通过合成提示诱导并行行为。

宽搜索：同时探索多个独立信息源，必须并行访问多源信息。

深搜索：多推理分支与延迟聚合，可并行探索不同推理路径。

长文档分析：超过100个输入文档的理解，可并行处理不同文档片段。

大规模文件下载：批量获取多样化资源，必须并发下载。

\Needspace{4\baselineskip}
\begin{enumerate}
\def\labelenumi{\arabic{enumi}.}
\setcounter{enumi}{4}
\tightlist
\item
  异步协程架构
\end{enumerate}

每个智能体任务作为独立异步协程运行，支持递归子任务调用（子智能体可进一步创建子智能体）。

\FloatBarrier
\Needspace{5\baselineskip}
\hypertarget{ux53c2ux8003ux6587ux732e-116}{%
\subsection{参考文献}\label{ux53c2ux8003ux6587ux732e-116}}

\vspace{0.25em}
\begin{itemize}
\tightlist
\item
  温睦宁、林江浩、张伟楠、俞勇. (2025). \href{https://haa.boyuai.com/}{《动手学大模型智能体》}. 人民邮电出版社. ISBN 978-7-115-68638-1.
\item
  Du, Y., Li, S., Torralba, A., Tenenbaum, J. B., \& Mordatch, I. (2023). \href{https://arxiv.org/abs/2305.14325}{Improving Factuality and Reasoning in Language Models through Multiagent Debate}. arXiv:2305.14325.
\item
  Wang, X., Wei, J., Schuurmans, D., et al.~(2023). \href{https://openreview.net/forum?id=1PL1NIMMrw}{Self-Consistency Improves Chain of Thought Reasoning in Language Models}. ICLR 2023.
\item
  Kimi Team. (2026). \href{https://arxiv.org/abs/2602.02276}{Kimi K2.5: Visual Agentic Intelligence}. arXiv:2602.02276.
\end{itemize}

\clearpage
